% from main.tex
%%%%%%%%%%  scan idx 379  |  folio 366  %%%%%%%%%%
sont linéairement dépendants sur \(K^\circ(S)\).

\(3^\circ\)) \(1,\ldots,\xi^r\) sont linéairement indépendants sur \(K^\circ(S)\).

\subsection*{1.2.}
Pour plus de commodité, nous allons rappeler ici quelques définitions et propriétés classiques des Modules gradués.

Soient \(S\) un schéma, \(B\) une \(\mathcal O_S\)-Algèbre graduée, \(M\) un \(B\)-Module gradué, Pour tout \(B\)-Module gradué \(N\), on désigne comme d'habitude par \(N(n)\) le \(B\)-Module gradué défini par :
\[
  (N(n))_k=N_{n+k}\quad\text{pour tout }k.
\]
On dit que \(M\) est un \(B\)-Module gradué localement libre (resp. localement libre de type fini) si tout point de \(S\) possède un voisinage ouvert \(U\) tel que \(M|U\) soit isomorphe, en tant que \(B\)-Module gradué, à une somme directe de Modules de la forme \(B(n)|U\) (resp. une somme directe finie).

On dit que \(M\) est un \(B\)-Module gradué de présentation finie si tout point de \(S\) possède un voisinage ouvert \(U\) tel qu'il existe une suite exacte \(P\longrightarrow Q\longrightarrow M|U\longrightarrow 0\), où \(P\) et \(Q\) sont des \((B|U)\)-Modules gradués libres de type fini, et les homomorphismes compatibles aux graduations.

\textbf{1.2.1.} Pour tout entier \(n\), on peut définir un homomorphisme canonique de \(B\)-Modules gradués :
\begin{equation}\tag{1.2.1}
  M_n\otimes_{\mathcal O_S} B(-n)\longrightarrow M,
\end{equation}
en envoyant \(x\otimes a\), où \(x\) est une section de \(M_n\), et \(a\) une section de \((B(-n))_k=B_{k-n}\) sur l'élément \(a.x\) de \(M_k\).

\textbf{1.2.2.} On dit qu'un \(B\)-Module gradué \(M\) est libre de type fini, et homogène de degré \(n\) s'il est isomorphe en tant que \(B\)-Module gradué à \((B(-n))^k\) pour un certain entier \(k\).

